
\section{Criterios locales de planitud}
\label{IV.5}

\begin{proposition}
\label{IV.5.1}
Sean~$A$ un anillo provisto de un ideal~$I$ y~$M$ un $A$-módulo.
Supongamos que
$$
\Tor_1^A(M,A/I^n)=0 \qquad \text{para $n>0$};
$$
entonces el homomorfismo canónico sobreyectivo
\begin{equation*}
\label{eq:IV.5.1.*}
\tag{$*$} {\gr_I^0(M)\otimes_{A/I}\gr_I(A)\to\gr_I(M)}
\end{equation*}
es un isomorfismo. El recíproco es cierto si~$I$ es nilpotente.
\end{proposition}

La hipótesis significa que los homomorfismos
$$
I^n\otimes_A M\to I^nM
$$
son isomorfismos, de donde se sigue inmediatamente que los homomorfismos
$$
I^n/I^{n+1}\otimes_A M\to I^nM/I^{n+1}M
$$
son isomorfismos. Recíprocamente, supongamos que se cumple esta condición
y que~$I$ es nilpotente; demostremos que~$\Tor_1^A(M,A/I^n)=0$ para todo
$n$. Esto es cierto para~$n$ grande; procedamos por inducción descendente
sobre~$n$, suponiéndolo demostrado para~$n+1$. Se tiene un diagrama
conmutativo
$$
\xymatrix{ & M\otimes I^{n+1}\ar[r]\ar[d] & M\otimes I^n\ar[r]
\ar[d] & M\otimes(I^n/I^{n+1})\ar[r]\ar[d] & 0 \\ 0\ar[r] &
MI^{n+1}\ar[r] & MI^n\ar[r] & MI^n/MI^{n+1}\ar[r] & 0 }
$$
cuyas filas son exactas.
\marginpar{96}
Por hipótesis, la última flecha vertical es un isomorfismo, y la hipótesis
de inducción significa asimismo que lo es la primera flecha vertical. Lo
mismo vale, por tanto, para la flecha vertical intermedia, lo que termina
la demostración.

Serre aisló la proposición siguiente durante el Seminario; permite
simplificaciones sustanciales en esta sección.
\begin{proposition}
\label{IV.5.2}
Sean~$A\to B$ un homomorfismo de anillos y~$M$ un $A$-módulo. Son
equivalentes las condiciones siguientes:
\begin{enumerate}
\item[(i)] Para todo $B$-módulo~$N$, se tiene $\Tor_1^A(M,N)=0$.
\item[(ii)] $\Tor_1^A(M,B)=0$ y $M_{(B)}=M\otimes_A B$ es plano
sobre~$B$.
\end{enumerate}
\end{proposition}

Se tiene un isomorfismo funtorial
$$
M\otimes_A N=(M\otimes_A B)\otimes_B N
$$
que expresa el primer miembro, considerado como funtor en~$M$, como
compuesto de los dos funtores $M\mto M\otimes_A B$ y
$P\mto P\otimes_B N$. Como el primero transforma módulos libres sobre~$A$
en módulos libres sobre~$B$, y por tanto proyectivos en proyectivos, se
tiene la sucesión espectral de los funtores compuestos
$$
\Tor_n^A(M,N)\From\Tor_p^B(\Tor_q^A(M,B),N),
$$
de donde resulta una sucesión exacta para los términos de grado bajo
$$
0\from\Tor_1^B(M\otimes_A B,N)\from\Tor_1^A(M,N)
\from\Tor_1^A(M,B)\ifthenelse{\boolean{orig}}{\otimes_A}{\otimes_B}N.
$$
Si se cumple~(i), entonces se concluye de esta sucesión exacta que
$\Tor_1^B(M\otimes_A B,N)=0$ para todo~$N$, es decir,
$M\otimes_A B$ es plano sobre~$B$, de donde~(ii). Si, recíprocamente, se
cumple~(ii), los términos que flanquean a~$\Tor_1^A(M,N)$ en la sucesión
exacta son nulos, y se obtiene~(i).

\begin{corollaire}
\label{IV.5.3}
Supongamos que~$B=A/I$; entonces las condiciones precedentes equivalen a
la siguiente:
\begin{enumerate}
\item[(iii)] $\Tor_1^A(M,N)=0$ para todo $A$-módulo~$N$ aniquilado por
una potencia de~$I$.
\end{enumerate}
\end{corollaire}

En efecto, (i) dice que esto vale cuando~$N$ está aniquilado por~$I$. Se
deduce~(iii) aplicando la hipótesis a los~$I^nN/I^{n+1}N$.

\begin{corollaire}
\label{IV.5.4}
Bajo las condiciones de~\Ref{IV.5.3}, las condiciones consideradas
implican que
\marginpar{97}
el homomorfismo funtorial
\begin{equation*}
\label{eq:IV.5.*}
\tag{$*$} {\gr_I^0(M)\otimes_{A/I}\gr_I(A)\to\gr_I(M)}
\end{equation*}
es un isomorfismo y que~$M\otimes_A A/I$ es plano sobre~$A/I$.
\end{corollaire}

Basta aplicar~(iii) y
\ifthenelse{\boolean{orig}}{\eqref{IV.5.1}}{\Ref{IV.5.1}}.
Utilizando el recíproco de~\Ref{IV.5.1} cuando~$I$ es nilpotente, se
obtiene:
\begin{corollaire}
\label{IV.5.5}
Sean~$A$ un anillo provisto de un ideal nilpotente~$I$ y~$M$ un
$A$-módulo. Son equivalentes las condiciones siguientes:
\begin{enumerate}
\item[(i)] $M$ es plano sobre~$A$.
\item[(ii)] $M\otimes_A A/I$ es plano sobre~$A/I$ y
$\Tor_1^A(M,A/I)=0$.
\item[(iii)] $M\otimes_A A/I$ es plano sobre~$A/I$ y el homomorfismo
canónico~\eqref{eq:IV.5.*} sobre los graduados asociados es un
isomorfismo.
\end{enumerate}
\end{corollaire}

Estas son, respectivamente, las condiciones (iii) y (ii) precedentes y
las condiciones del corolario~\Ref{IV.5.4}.

No supongamos ya que~$I$ sea nilpotente. Entonces, a priori, solo se
tienen en~\Ref{IV.5.5} las implicaciones
(i)$\To$(ii)$\To$(iii). Por otra parte, como la condición~(iii) permanece
estable al pasar al cociente por una potencia de~$I$, se ve, en virtud
de~\Ref{IV.5.5}, que implica
\begin{enumerate}
\item[(iv)] \emph{Para todo entero
\ifthenelse{\boolean{orig}}{$n$}{$n>0$},
$M\ifthenelse{\boolean{orig}}{\otimes}{\otimes_A}A/I^n$ es plano
sobre~$A/I^n$.}
\end{enumerate}
Nos proponemos dar condiciones bajo las cuales puede deducirse~(i), es
decir, que~$M$ es plano sobre~$A$. Afirmo que basta para ello que~$A$ sea
noetheriano y que~$M$ satisfaga la condición de finitud siguiente:
\emph{para todo módulo de tipo finito~$N$ sobre~$A$, $M\otimes_A N$ es
separado para la topología $I$-preádica}. (Bastaría verificarlo cuando~$N$
es un ideal de tipo finito de~$A$). En efecto, demostremos que, bajo estas
condiciones, si~$N'\to N$ es un monomorfismo de módulos de tipo finito,
$M\otimes_A N'\to M\otimes_A N$ es un monomorfismo. Basta mostrar que el
núcleo está contenido en todos los
$I^n(M\otimes_A N')=\Im(M\otimes_A I^nN'\to M\otimes_A N')$, o, lo que
equivale a lo mismo, en todos los
$\Im(M\otimes_A V'_n\to M\otimes_A N')=
\Ker(M\otimes_A N'\to M\otimes_A(N'/V'_n))$, donde~$V'_n$ recorre un
sistema fundamental numerable de entornos de~$0$ en~$N'$ (provisto de su
topología $I$-ádica). Por el teorema de Krull, la topología $I$-ádica
de~$N'$ es inducida por la de~$N$; se puede tomar, por tanto,
$V'_n=N'\cap I^nN$. Consideremos entonces el diagrama conmutativo
$$
\xymatrix{M\otimes_A N'\ar[r]\ar[d] & M\otimes_A(N'/V'_n)\ar[d]
\\ M\otimes_A N\ar[r] & M\otimes_A(N/I^nN)}.
$$
\marginpar{98}%
Como~$N'/V'_n$ y~$N/I^nN$ están aniquilados por~$I^n$, el segundo
homomorfismo vertical se identifica con el obtenido del homomorfismo
\emph{inyectivo} $N'/V'_n\to N/I^nN$ tensorizando sobre~$A/I^n$ con el
$(A/I^n)$-módulo \emph{plano} $M\otimes_A A/I^n$; es, por tanto,
\emph{inyectivo}. En consecuencia, el núcleo de
$M\otimes_A N'\to M\otimes_A N$ está contenido en el de
$M\otimes_A N'\to M\otimes_A(N'/V'_n)$, como se quería.

La condición de «finitud» considerada sobre~$M$ se cumple, en particular,
si~$M$ es un módulo de tipo finito sobre una $A$-álgebra noetheriana
\ifthenelse{\boolean{orig}}{,}{}~$B$ tal que~$IB$ esté contenido en el
radical de Jacobson de~$B$:
en efecto, entonces~$M\otimes_A N$ es un módulo de tipo finito sobre~$B$
para todo módulo de tipo finito~$N$ sobre~$A$, y por Krull es separado
para la topología $I$-ádica, que es su topología $(IB)$-ádica. Se obtiene
así:

\begin{theoreme}
\label{IV.5.6}
Sean~$A\to B$ un homomorfismo de anillos noetherianos, $I$ un ideal de~$A$
tal que~$IB$ esté contenido en el radical de Jacobson de~$B$ y~$M$ un
$B$-módulo de tipo finito. Son equivalentes las condiciones siguientes:
\begin{enumerate}
\item[(i)] $M$ es plano sobre~$A$.
\item[(ii)] $M\otimes_A A/I$ es plano sobre~$A/I$ y
$\Tor_1^A(M,A/I)=0$.
\item[(iii)] $M\otimes_A A/I$ es plano sobre~$A/I$ y el homomorfismo
canónico
$$
\gr_I^0(M)\otimes_{A/I}\gr_I(A)\to\gr_I(M)
$$
es un isomorfismo.
\item[(iv)] Para todo entero
\ifthenelse{\boolean{orig}}{$n$}{$n>0$}, $M\otimes_A A/I^n$ es plano
sobre~$A/I^n$.
\end{enumerate}
\end{theoreme}

Este resultado se aplica sobre todo cuando~$A$ y~$B$ son anillos
\emph{locales} noetherianos, $A\to B$ es un homomorfismo local e~$I$ es
un ideal de~$A$ contenido en su ideal maximal (y~\Ref{IV.5.6} puede
reducirse de inmediato a este caso). Un caso interesante es aquel en que
$A/I$ es un cuerpo, es decir, $I$ es maximal; entonces la condición de que
$M\otimes_A A/I$ sea plano sobre~$A/I$ es automática. Además, como los
$A/I^n$ son en tal caso anillos locales artinianos, la condición~(iv)
significa que los~$M\otimes_A A/I^n$ son \emph{libres} sobre~$A/I^n$.

\begin{corollaire}
\label{IV.5.7}
Sea
\marginpar{99}
$A\to B$ un homomorfismo local de anillos locales noetherianos y sea
$u\colon M'\to M$ un homomorfismo de $B$-módulos de tipo finito.
Supongamos que~$M$ es plano sobre~$A$. Son equivalentes las condiciones
siguientes:
\begin{enumerate}
\item[(i)] $u$ es inyectivo y~$\Coker u$ es plano sobre~$A$.
\item[(ii)] $u\otimes_A k\colon M'\otimes_A k\to M\otimes_A k$ es
inyectivo.
\end{enumerate}
(Aquí~$k$ denota el cuerpo residual de~$A$).
\end{corollaire}

(i)$\To$(ii) en virtud de~\Ref{IV.1.1}; demostremos el recíproco. En
primer lugar, $u$ es inyectivo, pues basta verificarlo en los graduados
asociados, donde resulta de un cuadrado conmutativo que el lector
escribirá. Sea~$M''$ su conúcleo; se tiene, por tanto, una sucesión exacta
$$
0\to M'\to M\to M''\to0,
$$
de donde la sucesión exacta de los~$\Tor$, teniendo en cuenta la
hipótesis~(ii) y~$\Tor_1^A(M,k)=0$, da la relación
$\Tor_1^A(M'',k)=0$; por el teorema~\Ref{IV.5.6}, $M''$ es plano
sobre~$A$.

\begin{corollaire}
\label{IV.5.8}
Bajo las condiciones de~\Ref{IV.5.6}, sea~$J$ un ideal de~$B$ que
contiene a~$IB$ y está contenido en el radical de Jacobson. Sean
$\hat A$ la completación
\ifthenelse{\boolean{orig}}{$I$-ádicas}{$I$-ádica} de~$A$ y~$\hat B$
y~$\hat M$ las completaciones $J$-ádicas de~$B$ y~$M$. Para que~$M$ sea
plano sobre~$A$, es necesario y suficiente que
\ifthenelse{\boolean{orig}}{$M$}{$\hat M$} sea plano sobre~$\hat A$.
\end{corollaire}

(N.B. la suficiencia resultaría ya fácilmente de~\Ref{IV.3.2}). Se
utiliza el criterio~(iii) de~\Ref{IV.5.6} en la situación~$(A,B,I,M)$ y
en la situación~$(\hat A,\hat B,I\hat A,\hat M)$. Se comprueba que las
condiciones obtenidas en ambos casos son equivalentes gracias
a~\Ref{IV.3.2}.

\begin{corollaire}
\label{IV.5.9}
Sean~$A\to B\to C$ homomorfismos locales de anillos locales
noetherianos y~$M$ un $C$-módulo de tipo finito (N.B. $C$ interviene
únicamente para poder imponer una condición de finitud a~$M$).
Supongamos que~$B$ es plano sobre~$A$. Sea~$k$ el cuerpo residual de~$A$.
Son equivalentes las condiciones siguientes:
\begin{enumerate}
\item[(i)] $M$ es plano sobre~$B$.
\item[(ii)] $M$ es plano sobre~$A$ y~$M\otimes_A k$ es plano
sobre~$B\otimes_A k$.
\end{enumerate}
\end{corollaire}

La implicación (i)$\To$(ii) es trivial; demostremos (ii)$\To$(i). Se
aplica el criterio~(iii) de~\Ref{IV.5.6} a~$(B,C,\goth{m}B=I,M)$. Como
$M\otimes_B(B/I)=M\otimes_B(B\otimes_A k)=M\otimes_A k$, la primera
condición de este criterio significa precisamente que~$M\otimes_A k$ es
plano sobre~$B\otimes_A k$, perfecto.
\marginpar{100}
La segunda condición del criterio se cumple porque~$M$ es plano sobre~$A$
y~$B$ es plano sobre~$A$, por una fórmula de asociatividad del producto
tensorial. --- Naturalmente, remitiéndose a~\Ref{IV.5.5} en lugar de
a~\Ref{IV.5.6}, se obtiene un enunciado análogo sin condición noetheriana
ni de finitud, si se supone, en cambio, que el ideal~$\goth{m}$ de~$A$ es
nilpotente. (El hecho de que se haya tomado~$\goth{m}$ maximal no ha
intervenido; pero, en cierto sentido, el caso «$\goth{m}$ maximal» es «el
mejor posible»).
